\section{Conclusion}

We present an empirical study of OPD and OPSD for LLMs, showing that the effectiveness depends on task structure and the nature of privileged information. OPSD is ineffective in our tested mathematical reasoning settings, where privileged information is largely instance-specific, but works well for system prompt internalization and alignment tasks where privileged information reflects shared latent rules. We identify three failure mechanisms: teacher-student distribution mismatch, biased TopK reverse-KL gradients, and the limitation of OPSD in learning only the marginal distribution over instance-specific privileged information. We further show that stop-gradient TopK surrogate loss, RLVR teacher improvement, and SFT stabilization can improve practical training stability. Our findings are based on a limited set of model families and model scales, and larger-scale settings may exhibit different behavior. Future work may explore iterative self-improvement pipelines that combine SFT, RL, and OPD: SFT provides a stable initialization, RL improves task-specific behavior, and OPD distills the improved on-policy behavior back into the student.